% LaTeX fragment for the parent manuscript.  This sharpens the incoming
% general k+2 bound; the delivered production checker remains threshold three.

\section{Exact universal cutoffs for connected tiles}
\label{sec:exact-cutoffs}

The incoming article established the universal connected-tile cutoff $k+2$
and improved it sharply to four at threshold three.  Its further questions
included the optimal cutoffs at other thresholds.  For the same unit-edge
model with two exact monochromatic-edge marks, that question has the following
complete answer.

\begin{theorem}[Exact connected-tile stabilization threshold]
\label{thm:exact-cutoff}
For a nonnegative integer $k$, put
\[
 b_k=\begin{cases}
 1,&k=0,\\
 3,&k=1,\\
 4,&k=2,\\
 k+1,&k\geq3.
 \end{cases}
\]
For every connected tile $T$ and every positive integer $N$,
\[
 \mathsf S_k(T^{\circ N})
   =\mathsf S_k(T^{\circ\min(N,b_k)}).
\]
No smaller cutoff works uniformly for all connected tiles in this model.
\end{theorem}

\begin{proof}[Upper bound]
A connected tile has at least one edge because its nonempty input and output
interfaces are disjoint.  A zero-cost coloring of any power is monochromatic,
has one exact edge mark, and prescribes that same color on both interfaces.
This proves the assertion at $k=0$.  For $k=1,2$, start with $N>k+2$
copies and a coloring of cost at most $k$.  Retain all at most $k$
positive-cost copies and at most one zero-cost copy of each color that
appeared.  The latter copies are monochromatic, so their deletion does not
change seam colors, and retaining one of each occurring color preserves
the exact marks.  At most $k+2$ copies remain, including at least one
zero-cost copy; duplicate that copy to reach length $k+2$.  Conversely,
any coloring of $k+2$ copies of cost at most $k$ has a zero-cost copy
that can be duplicated to obtain every larger length.  This proves the
two remaining small-threshold upper bounds.

Suppose $k\geq3$ and $N>k+1$.  Fix a coloring of $T^{\circ N}$ with total
cost at most $k$.  At most $k$ tile occurrences have positive cost.  Every
other occurrence has zero cost, is monochromatic, and has identical constant
input and output colors.  Consequently it can be deleted from its serial
position without creating a coloring disagreement at the new seam.
Keep all positive-cost occurrences and at most one zero-cost occurrence of
each color that appeared.  This preserves boundary colors, cost, and exact
marks.  At least one zero-cost occurrence remains because $N>k$.

Let $p$ be the number of retained positive-cost occurrences and $z\leq2$
the number of retained zero-cost occurrences.  If $p+z\leq k+1$, duplicate
one retained zero-cost occurrence until the length is exactly $k+1$.
The resulting coloring preserves all the data and its cost.

The only exceptional case is $p=k$ and $z=2$.  Each positive-cost occurrence
then has cost exactly one.  If some positive-cost occurrence supplies a
monochromatic-edge mark $a\in\{0,1\}$, delete the retained zero-cost
occurrence of color $a$.  That mark is still supplied by the positive-cost
occurrence, and the other retained zero-cost occurrence supplies the other
mark.  The new length is exactly $k+1$ and no required data change.

Otherwise every positive-cost occurrence has no monochromatic edge.  All
edges of the tile are therefore bichromatic in that occurrence.  Its cost
is one, so the tile has exactly one edge.  Connectedness and the disjoint
nonempty interfaces force $T$ to be the single-edge tile with one input and
one output.  The two retained zero-cost occurrences imply that the required
exact mark set is $\{0,1\}$.  For either prescribed pair of endpoint colors,
there is a path of length at most four with both marks and cut cost at most
two: use colors $0,0,1,1,0$ for equal zero endpoints, its complement for
equal one endpoints, $0,0,1,1$ for distinct endpoints, or its reversal.
These paths contain a monochromatic edge and may be padded to length $k+1$
by duplicating it.  The cost is at most two, hence at most $k$, and the
boundary colors and exact mark set are the required ones.

We have proved that every feasible source coloring of cost at most $k$
gives a coloring at length $k+1$ of no greater cost with the same boundary
colors and exact marks.  Conversely, a coloring of $k+1$ copies with cost
at most $k$ contains a zero-cost occurrence.  Duplicating it produces a
coloring at any larger length with exactly the same cost, boundary colors,
and marks.  These two directions establish equality of all minimum costs
up to the threshold and equality of all saturated entries.

The argument preserves minimum signatures; it does not claim that every
individual source coloring can be shortened while preserving its exact cost.
The single-edge exceptional case is the only step where its cost may decrease.
\end{proof}

\begin{proof}[Sharpness]
At threshold zero, the smallest possible positive cutoff is one.
For threshold one, take the single-edge tile with distinct endpoint colors
and require both marks.  A path of three edges with colors $0,0,1,1$ has
cost one and both marks.  No path of at most two edges has these properties:
it needs one edge of each monochromatic color and at least one crossing edge.
Thus $b_1\geq3$.

For threshold two, prescribe both endpoints zero and again require both
marks.  Four edges with colors $0,0,1,1,0$ suffice.  Any such coloring needs
an edge of each monochromatic color and at least two crossing edges, so at
least four edges are necessary.  Hence $b_2\geq4$.

Now let $k\geq3$ and put $B=k+1$.  Form a connected tile $T_k$ of width
$2B$ with input vertices $\ell_1,\ldots,\ell_{2B}$, output vertices
$r_1,\ldots,r_{2B}$, and edges
\[
 \ell_i r_i \quad(1\leq i\leq2B),
 \qquad
 \ell_i\ell_{i+1}\quad(1\leq i<2B).
\]
Thus the inputs form a path and every output is attached to its corresponding
input by one edge.  Prescribe color one on the first $B$ input and output
coordinates, and color zero on all the remaining coordinates.

In $T_k^{\circ k}$, propagate these colors unchanged along each vertical
strand.  Exactly the input-path edge joining coordinates $B$ and $B+1$
is cut in each copy.  The total cost is $k$, and the matching edges supply
both exact marks.

In contrast, suppose that $T_k^{\circ(k+1)}$ had any coloring of cost at
most $k$ with those boundary colors.  Some connected tile occurrence would
have zero cost and hence be uniformly colored $a$.  For each of the $B$
oppositely colored original input coordinates, follow its matching-edge
strand from the exposed input to the corresponding input of that occurrence.
If the occurrence is the first one, its monochromatic input condition already
contradicts the prescribed colors.  Otherwise these are $B$ pairwise
edge-disjoint paths, each with oppositely colored endpoints and hence each
containing a cut edge.  The total cut cost is at least $B=k+1$, a contradiction.

Therefore the chosen marked-signature entry is at most $k$ at exponent $k$
and is saturated at exponent $k+1$.  If a universal cutoff were at most $k$,
both exponents would reduce to the same capped exponent and would have equal
signatures.  This contradiction proves $b_k\geq k+1$.
\end{proof}

\paragraph{Consequences and implementation scope.}
Every routed or reblocked kernel theorem can use $b_k$ in place of $a_k$.
In particular the routed explicit-size coefficient $a_k+1$ becomes $k+2$
for all $k\geq3$.  The threshold-three value remains four, so the production
cubic-pattern checker and its source witness format require no change.
The supplied general-threshold producer uses the exact value $b_k$ in
\texttt{connected\_cutoff}; both aligned and reblocked construction use it.
The certificate solver remains the threshold-three cubic-pattern query.
The comb examples establish sharpness in the
full connected unit-edge graph-tile class, without a claim of optimality for
more restrictive cubic or geometrically realizable subclasses.
